\honest{Cached Thoughts}{Eliezer Yudkowsky}{2007}

Neurons fire at most about two hundred times a second, I begin, so a mental operation can take at most about a hundred steps in a row. A programmer stuck with processors that slow would store the results of past work and look them up, instead of working them out again. This is called caching. ``It's a good guess,'' I say, that most human thinking consists of looking things up. After this I use only one phrase from it, ``the brain completes the pattern.''

My example is a story I once read and cannot find again. A know-it-all neighbor told a man that the way to remove a chimney is to knock out the fireplace and let the bricks drop, one level at a time. Years later the man tried it, and it did not go well. \nb{This is a man who acted on a bad source without checking it.} From here on, a ``cached thought'' means believing what someone told you.

I suspect the man would have been more critical of the idea if it had been his own. \nb{An experiment by Trouche, Mercier and colleagues suggests the reverse: people judge other people's arguments more harshly than their own. Yudkowsky's own posts of the past two weeks, on the clever arguer and on rationalization, are about people defending their own conclusions.}

Then I concede that nobody can think fast enough to work out everything alone; raised by wolves, I would hardly be human. \nb{So relying on other people's conclusions is how knowledge works, and the real question is which conclusions to trust.} My only answer is to distrust thoughts ``not invented by critical thinkers.''

Next come my examples of cached thoughts. The first, I say, is repeated even by people who aspire to critical thinking: ``You can't prove or disprove a religion by factual evidence.'' ``Death gives meaning to life.'' ``Maybe the human species doesn't deserve to survive,'' which decent people say when I raise the risk of human extinction. ``Love isn't rational.'' Each one is an opinion I argue against elsewhere. \nb{The first had a well-known, argued defense by the scientist Stephen Jay Gould.} I do not mention him; a post I link in a footnote argues against the idea. Here I add that a few centuries ago the claim would have got you burned at the stake, and I write a prayer for a mother whose daughter has cancer, in which she recites the skeptic's philosophy to God.

For ``Love isn't rational,'' I ask how you would examine the idea if it were new to you. ``I know what I would say,'' I add, with a link to my answer.

At the end I tell you what the post is for. Next time ``you hear someone unhesitatingly repeating a meme you think is silly or false,'' you will think, ``Cached thoughts.'' Earlier I asked what patterns are being completed ``inside your mind.'' But the use I name at the end is for views you already reject, held by other people. I then ask whether my own idea is true, and tell you to think.
